\section{Gauge Symmetry and Why It Forbids a Photon Mass}

The theoretical foundation of electromagnetism is the $U(1)$ gauge symmetry. 
This symmetry is not an aesthetic choice but the structural principle that 
determines the form of Maxwell’s equations, guarantees charge conservation, 
and fixes the number of physical degrees of freedom of the electromagnetic 
field. A key consequence of this symmetry is that a photon mass term is 
incompatible with an unbroken $U(1)$ gauge theory.

\subsection*{Gauge symmetry and the structure of the Lagrangian}
The Lagrangian of free electromagnetism is
\[
\mathcal{L}_{\mathrm{EM}} = -\frac{1}{4}F_{\mu\nu}F^{\mu\nu},
\]
where the field strength tensor is defined as 
$F_{\mu\nu} = \partial_\mu A_\nu - \partial_\nu A_\mu$.
This theory is invariant under local gauge transformations,
\[
A_\mu(x) \;\rightarrow\; A_\mu(x) + \partial_\mu \Lambda(x).
\]
A photon mass term of the form
\[
\mathcal{L}_{\mathrm{mass}} = \frac{1}{2}m_\gamma^2 A_\mu A^\mu
\]
is \emph{not} invariant under this transformation, because the additional 
$\partial_\mu \Lambda$ contribution does not vanish. Introducing such a 
term explicitly breaks gauge symmetry.

\subsection*{Physical consequences of breaking gauge symmetry}
Breaking gauge symmetry is not a small modification—it fundamentally alters 
the structure of the theory:

\begin{itemize}
	\item charge conservation is no longer guaranteed;
	\item Maxwell’s equations acquire additional terms and lose their 
	elegant differential structure;
	\item the photon gains an additional degree of freedom: a massive 
	spin-1 particle has three polarization states instead of two;
	\item the electromagnetic interaction acquires a finite range, 
	analogous to a Yukawa potential.
\end{itemize}

These consequences are not abstract mathematical curiosities; they would 
lead to clear experimental signatures, many of which are tightly constrained 
by terrestrial and astrophysical observations.

\subsection*{Why an unbroken $U(1)$ symmetry enforces a massless photon}
Modern field theory permits mass terms only if they are compatible with the 
symmetry structure or generated through spontaneous symmetry breaking. 
In the electroweak sector, the $W^\pm$ and $Z$ bosons acquire mass through 
the Higgs mechanism. However, the electromagnetic $U(1)$ symmetry remains 
unbroken, and therefore the photon cannot obtain a mass in this framework.

The masslessness of the photon is thus not an arbitrary assumption but a 
direct consequence of the underlying gauge symmetry. Any deviation from 
$m_\gamma = 0$ would require physics beyond the Standard Model—specifically, 
a mechanism capable of breaking or modifying the electromagnetic gauge group.

\subsection*{Position within the overall argument}
The conclusion of this chapter is straightforward: a photon mass term is 
incompatible with the gauge structure of electromagnetism. The next chapter 
examines what would happen if we nevertheless introduced such a term. The 
resulting Proca theory provides a concrete framework to explore the physical 
implications of a massive photon and to connect these theoretical predictions 
with empirical constraints.

Gauge symmetry therefore provides the conceptual foundation for understanding 
why the photon is expected to be massless—and why even a tiny deviation from 
this prediction would point to new and unexplored physics.

